\documentclass[UTF8,11pt]{ctexart}
\usepackage{amsmath,amssymb,mathtools}
\usepackage{geometry}
\usepackage{bm}
\geometry{a4paper,margin=2.5cm}
\allowdisplaybreaks

\title{STEP1：单支路冲击传播\\
第二版事件定义与第一版对称公式的统一}
\author{}
\date{}

\begin{document}
\maketitle

\section{目的与统一口径}

本文保留两类互补的表达方式：

\begin{enumerate}
    \item 用明确的事件定义区分“生产节点自身何时停产/恢复”与
    “该冲击何时真正被某个具体下游观察到”；
    \item 在单一停产窗口下，保留原有版本高度对称的
    stop/recovery 递推与累加公式；
    \item 在多窗口与一般 DAG 下，再用累计缺口函数 \(D(t)\)
    作为上述公式的统一推广。
\end{enumerate}

核心原则是
\[
\boxed{
\text{先定义物理事件，再给单窗口闭式公式，最后推广为 \(D(t)\) 算子。}
}
\]

考虑一条直接生产关系分支
\[
P_p\longrightarrow M_1\longrightarrow\cdots
\longrightarrow M_K\longrightarrow P_q.
\]

\section{两类事件必须严格区分}

\subsection{生产节点自身的 production-start event}

记 \(u_p(t)\) 为生产节点 \(P_p\) 在时刻 \(t\) 实际启动的生产活动强度，
初始化正常强度为
\[
u_p^0\equiv u_p(0).
\]

若 \(P_p\) 在一个单一完全停产窗口内停止启动新生产，定义
\[
\boxed{
I_p^{act}=[\sigma_p,\rho_p),
}
\]
其中
\[
\boxed{
\sigma_p
=
\text{\(P_p\) 自身停止启动新生产的时刻},
}
\]
\[
\boxed{
\rho_p
=
\text{\(P_p\) 自身重新以正常强度启动生产的时刻}.
}
\]

定义该内部停产窗口长度
\[
\boxed{
\kappa_p=\rho_p-\sigma_p.
}
\]

这里的 \(\sigma_p,\rho_p\) 只描述
\[
\boxed{\text{\(P_p\) 自身的 production-start plane}}
\]
上的事件，不表示任何具体下游已经感受到断供。

\subsection{具体下游观察到的 stop/recovery}

对具体直接下游 \(P_q\)，定义
\[
\boxed{
t_{p\to q}^{\mathrm{stop}}
=
\text{由 \(P_p\) 的供给缺口导致
\(P_q\) 首次无法继续获得正常所需投入的时刻},
}
\]
以及
\[
\boxed{
t_{p\to q}^{\mathrm{rec}}
=
\text{\(P_q\) 在发生有效断供后，
第一次重新收到由 \(P_p\) 恢复生产产生的正常供给的时刻}.
}
\]

因此必须区分
\[
\boxed{
(\sigma_p,\rho_p)
\neq
(t_{p\to q}^{\mathrm{stop}},t_{p\to q}^{\mathrm{rec}}).
}
\]

前者是
\[
\text{node-local event},
\]
后者是
\[
\text{downstream-visible event}.
\]

若 \(P_p\) 同时连接多个下游，则一般有
\[
t_{p\to q_1}^{\mathrm{stop}}
\neq
t_{p\to q_2}^{\mathrm{stop}},
\qquad
t_{p\to q_1}^{\mathrm{rec}}
\neq
t_{p\to q_2}^{\mathrm{rec}},
\]
因为不同输出分支可以具有不同的 delay 与 inventory buffer。

\section{两个基本传播量}

对直接分支 \(P_p\to P_q\)，记运输边集合和库存节点集合分别为
\[
\mathcal E_{p\to q},
\qquad
\mathcal M_{p\to q}.
\]

\subsection{总时间平移}

定义
\[
\boxed{
\Theta_{p\to q}
=
\tau_p^{prod}
+
\sum_{e\in\mathcal E_{p\to q}}
\tau_e^{trans}.
}
\tag{1}
\]

这里把上游 \(P_p\) 自身的 production delay 明确归入
\(P_p\to P_q\) 的 downstream propagation。

因此 \(\Theta_{p\to q}\) 只负责
\[
\boxed{\text{时间平移}}
\]
而不会改变缺口长度。

\subsection{总库存 coverage}

对每个库存节点 \(m\in\mathcal M_{p\to q}\)，
其初始化正常流出率记为
\[
f_m^{out,0}>0.
\]

定义库存 coverage
\[
c_m=\frac{q_m(0)}{f_m^{out,0}},
\]
于是整条分支的总库存覆盖时间为
\[
\boxed{
C_{p\to q}
=
\sum_{m\in\mathcal M_{p\to q}}
\frac{q_m(0)}{f_m^{out,0}}.
}
\tag{2}
\]

物理上：
\[
\boxed{
\text{delay 决定冲击何时到达，inventory 决定冲击还能剩多长。}
}
\]

\section{STEP1 主结果：单一停产窗口的对称传播公式}

设 \(P_p\) 自身的 production-start gap 为
\[
[\sigma_p,\rho_p),
\qquad
\kappa_p=\rho_p-\sigma_p.
\]

定义候选停供前沿
\[
\boxed{
s_{p\to q}
=
\sigma_p+\Theta_{p\to q}+C_{p\to q},
}
\tag{3}
\]
以及恢复产出到达前沿
\[
\boxed{
r_{p\to q}
=
\rho_p+\Theta_{p\to q}.
}
\tag{4}
\]

如果
\[
s_{p\to q}<r_{p\to q},
\]
等价于
\[
\boxed{
\kappa_p>C_{p\to q},
}
\tag{5}
\]
则此次 \(P_p\) 的内部停产真正穿透该分支全部库存并影响到 \(P_q\)。

此时 STEP1 的核心公式为
\[
\boxed{
\begin{aligned}
t_{p\to q}^{\mathrm{stop}}
&=
\sigma_p
+\Theta_{p\to q}
+C_{p\to q},
\\[1mm]
t_{p\to q}^{\mathrm{rec}}
&=
\rho_p
+\Theta_{p\to q}.
\end{aligned}
}
\tag{6}
\]

这保留了原有版本最直接的对称结构：
\[
\boxed{
\begin{aligned}
\text{stop}
&=
\text{upstream stop}
+\text{delay}
+\text{inventory coverage},
\\
\text{recovery}
&=
\text{upstream recovery}
+\text{delay}.
\end{aligned}
}
\tag{7}
\]

两式相减得到
\[
\boxed{
t_{p\to q}^{\mathrm{rec}}
-
t_{p\to q}^{\mathrm{stop}}
=
(\rho_p-\sigma_p)-C_{p\to q}.
}
\tag{8}
\]

因此有效断供持续时间为
\[
\boxed{
\Delta_{p\to q}
=
\left[
(\rho_p-\sigma_p)-C_{p\to q}
\right]_+.
}
\tag{9}
\]

如果
\[
s_{p\to q}\ge r_{p\to q},
\]
则
\[
\boxed{
\Delta_{p\to q}=0.
}
\tag{10}
\]

此时 \(P_p\) 自身虽然确实停产，
但恢复产出前沿在库存完全耗尽之前已经追上缺口前沿，
所以 \(P_q\) 从未经历一次真正的有效断供。
此时 \(r_{p\to q}\) 可以称为“恢复产出到达前沿”，
但不应说 \(P_q\) 发生了“先停后恢复”。

\section{标准 \(P-M-M-P\) 结构}

考虑
\[
P_p\rightarrow M_1\rightarrow M_2\rightarrow P_q.
\]

定义
\[
\boxed{
\begin{aligned}
\Theta_{p\to q}
={}&
\tau_p^{prod}
+\tau_{P_p\to M_1}^{trans}
+\tau_{M_1\to M_2}^{trans}
+\tau_{M_2\to P_q}^{trans},
\\[1mm]
C_{p\to q}
={}&
\frac{q_{M_1}(0)}
     {f_{M_1\to M_2}^{0}}
+
\frac{q_{M_2}(0)}
     {f_{M_2\to P_q}^{0}}.
\end{aligned}
}
\tag{11}
\]

若单一停产窗口能够穿透全部库存，则
\[
\boxed{
\begin{aligned}
t_{p\to q}^{\mathrm{stop}}
={}&
\sigma_p
+\tau_p^{prod}
+\tau_{P_p\to M_1}^{trans}
+\frac{q_{M_1}(0)}{f_{M_1\to M_2}^{0}}
\\
&+
\tau_{M_1\to M_2}^{trans}
+\frac{q_{M_2}(0)}{f_{M_2\to P_q}^{0}}
+\tau_{M_2\to P_q}^{trans},
\\[2mm]
t_{p\to q}^{\mathrm{rec}}
={}&
\rho_p
+\tau_p^{prod}
+\tau_{P_p\to M_1}^{trans}
+\tau_{M_1\to M_2}^{trans}
+\tau_{M_2\to P_q}^{trans}.
\end{aligned}
}
\tag{12}
\]

因此 production delay 与 transport delay
在 stop/recovery 两个前沿中完全对称出现；
库存 coverage 只进入 stop 前沿。

\section{单源直链：恢复原有版本的递推与累加格式}

考虑
\[
P_0
\rightarrow\mathcal B_1
\rightarrow P_1
\rightarrow\mathcal B_2
\rightarrow P_2
\rightarrow\cdots
\rightarrow\mathcal B_i
\rightarrow P_i,
\]
其中 \(\mathcal B_j\) 表示
\(P_{j-1}\) 到 \(P_j\) 之间不含其他生产节点的一条物料支路。

定义
\[
\Theta_j
\equiv
\Theta_{P_{j-1}\to P_j},
\qquad
C_j
\equiv
C_{P_{j-1}\to P_j}.
\]

设源头 \(P_0\) 自身发生一个单一 production-start gap
\[
[\sigma_0,\rho_0).
\]

在单输入直链中，只要冲击尚未被库存完全吸收，
\(P_{j-1}\) 对 \(P_j\) 的 downstream-visible shortage
就直接决定 \(P_j\) 自身的 production-start gap。
因此有递推式
\[
\boxed{
\begin{aligned}
\sigma_j
&=
\sigma_{j-1}
+\Theta_j
+C_j,
\\[1mm]
\rho_j
&=
\rho_{j-1}
+\Theta_j.
\end{aligned}
}
\tag{13}
\]

逐段累加得到
\[
\boxed{
\begin{aligned}
\sigma_i
&=
\sigma_0
+\sum_{j=1}^{i}\Theta_j
+\sum_{j=1}^{i}C_j,
\\[1mm]
\rho_i
&=
\rho_0
+\sum_{j=1}^{i}\Theta_j.
\end{aligned}
}
\tag{14}
\]

因此
\[
\boxed{
\rho_i-\sigma_i
=
(\rho_0-\sigma_0)
-
\sum_{j=1}^{i}C_j.
}
\tag{15}
\]

若旧版设定为
\[
\sigma_0=0,
\qquad
\rho_0=k,
\]
则立刻得到
\[
\boxed{
\begin{aligned}
\sigma_i
&=
\sum_{j=1}^{i}(\Theta_j+C_j),
\\[1mm]
\rho_i
&=
k+\sum_{j=1}^{i}\Theta_j,
\end{aligned}
}
\tag{16}
\]
以及
\[
\boxed{
\rho_i-\sigma_i
=
k-\sum_{j=1}^{i}C_j.
}
\tag{17}
\]

若
\[
k\le \sum_{j=1}^{i}C_j,
\]
则该冲击在到达 \(P_i\) 之前已经被累计库存 coverage 完全吸收，
因此 \(P_i\) 不形成有效 production-start gap。

\section{从原有累加公式推广到累计缺口函数 \(D(t)\)}

原有对称公式适合描述
\[
\boxed{\text{单一 stop/recovery 窗口}.}
\]

为了处理多个不连续窗口、汇合后再分叉以及一般 DAG，
定义 \(P_p\) 的内部累计 production-start 缺口
\[
\boxed{
D_p^{act}(t)
=
\text{截至 \(t\)，\(P_p\) 相对于正常状态累计少启动的等效正常生产时间}.
}
\tag{18}
\]

在完全停产---完全恢复 baseline 下，
\[
\dot D_p^{act}(t)\in\{0,1\}
\]
几乎处处成立，并且
\[
\boxed{
u_p(t)
=
u_p^0\left[1-\dot D_p^{act}(t)\right].
}
\tag{19}
\]

对于单一内部停产窗口
\[
[\sigma_p,\rho_p),
\]
有
\[
\boxed{
D_p^{act}(t)
=
\min\left\{
(t-\sigma_p)_+,
\rho_p-\sigma_p
\right\}.
}
\tag{20}
\]

约定 \(s<0\) 时
\[
D_p^{act}(s)=0.
\]

则任意窗口情况下的 STEP1 传播算子定义为
\[
\boxed{
D_{p\to q}(t)
=
\mathcal T_{p\to q}[D_p^{act}](t)
=
\left[
D_p^{act}(t-\Theta_{p\to q})
-
C_{p\to q}
\right]_+.
}
\tag{21}
\]

这里
\[
(x)_+\equiv\max\{x,0\}.
\]

其基本构件为
\[
\boxed{
\text{Delay}:D(t)\mapsto D(t-\tau),
}
\tag{22}
\]
以及
\[
\boxed{
\text{Inventory}:D(t)\mapsto
\left[D(t)-\frac{q(0)}{f^0}\right]_+.
}
\tag{23}
\]

由于
\[
[[x-c_1]_+-c_2]_+
=
[x-c_1-c_2]_+,
\]
多个串联库存 coverage 可以直接累加。

\section{\(D(t)\) 算子如何退化回原有对称公式}

将单一窗口式 \((20)\) 代入 \((21)\)：
\[
D_{p\to q}(t)
=
\left[
\min\left\{
(t-\Theta_{p\to q}-\sigma_p)_+,
\rho_p-\sigma_p
\right\}
-
C_{p\to q}
\right]_+.
\tag{24}
\]

若
\[
\rho_p-\sigma_p>C_{p\to q},
\]
则 \(D_{p\to q}(t)\) 开始增长的时刻正是
\[
\boxed{
t_{p\to q}^{\mathrm{stop}}
=
\sigma_p+\Theta_{p\to q}+C_{p\to q},
}
\]
而停止增长的时刻正是
\[
\boxed{
t_{p\to q}^{\mathrm{rec}}
=
\rho_p+\Theta_{p\to q}.
}
\]

因此
\[
\boxed{
\text{\(D(t)\) 不是另一套模型，
而是原有 stop/recovery 累加公式的多窗口推广。}
}
\tag{25}
\]

\section{算子复合仍然就是原有的“逐段累加”}

定义
\[
\mathcal T_{\Theta,C}[D](t)
=
[D(t-\Theta)-C]_+.
\]

经过两段支路：
\[
D_1(t)
=
\mathcal T_{\Theta_1,C_1}[D_0](t),
\]
\[
D_2(t)
=
\mathcal T_{\Theta_2,C_2}[D_1](t).
\]

利用
\[
[[x-C_1]_+-C_2]_+
=
[x-C_1-C_2]_+,
\]
得到
\[
\boxed{
D_2(t)
=
\left[
D_0(t-\Theta_1-\Theta_2)
-(C_1+C_2)
\right]_+.
}
\tag{26}
\]

因此沿单源直链连续复合得到
\[
\boxed{
D_i^{act}(t)
=
\left[
D_0^{act}
\left(
t-\sum_{j=1}^{i}\Theta_j
\right)
-
\sum_{j=1}^{i}C_j
\right]_+.
}
\tag{27}
\]

所以新算子在没有汇合的直链上，
与旧版“delay 逐段累加、inventory coverage 逐段累加”的计算完全一致。

\section{与原有“第 \(i\) 个厂第 \(k\) 天生产状况”的对应}

\(D_i^{act}(t)\) 是原有逐日生产状态的连续时间累计表示。

在二元 baseline 下，
\[
\boxed{
\frac{u_i(t)}{u_i^0}
=
1-\dot D_i^{act}(t).
}
\tag{28}
\]

因此
\[
\dot D_i^{act}(t)=0
\quad\Longleftrightarrow\quad
P_i\text{ 正常生产},
\]
\[
\dot D_i^{act}(t)=1
\quad\Longleftrightarrow\quad
P_i\text{ 完全停产}.
\]

若把第 \(k\) 天定义为
\[
[k,k+1),
\]
并且该日内部没有额外的日内切换，则可写
\[
\boxed{
x_i[k]
=
1-
\left(
D_i^{act}(k+1)-D_i^{act}(k)
\right),
}
\tag{29}
\]
其中
\[
x_i[k]\in\{0,1\}.
\]

于是
\[
\boxed{
x_i[k]=1
\Longleftrightarrow
\text{第 \(i\) 个厂第 \(k\) 天生产},
}
\]
\[
\boxed{
x_i[k]=0
\Longleftrightarrow
\text{第 \(i\) 个厂第 \(k\) 天停产}.
}
\]

若一天内部允许发生 stop/recovery 切换，则
\[
D_i^{act}(k+1)-D_i^{act}(k)\in[0,1]
\]
自然表示该日损失的正常生产时间比例。

\section{多窗口情况下如何定义 stop/recovery}

如果 \(P_p\) 经历多个不连续 production-start gaps，
同一份下游库存不能针对每个窗口重新初始化为 \(q(0)\)。

因此不能对每个窗口分别套用一次
\[
C_{p\to q}
\]
再把得到的区间简单取并集。

正确做法是先计算完整的
\[
D_{p\to q}(t)
=
[D_p^{act}(t-\Theta_{p\to q})-C_{p\to q}]_+.
\]

在完全停产---完全恢复 baseline 下，
所有满足
\[
\dot D_{p\to q}(t)=1
\]
的最大连通区间定义为
\[
\boxed{
[t_{p\to q}^{\mathrm{stop},\ell},
t_{p\to q}^{\mathrm{rec},\ell}),
\qquad
\ell=1,\ldots,L_{p\to q}.
}
\tag{30}
\]

所以一般情形下，
\[
\boxed{
D_{p\to q}(t)
\text{ 是主接口，}
\qquad
(t^{\mathrm{stop}},t^{\mathrm{rec}})
\text{ 是从 \(D\) 中读取的事件标记。}
}
\]

\section{正常流率与固定比例}

若 \(P_p\) 的某种产出 \(Y\) 的产出系数为
\[
b_{pY}>0,
\]
并以固定比例
\[
\rho_{p\to q}>0
\]
分配到通往 \(P_q\) 的分支，且输出边 capacity 不绑定，
则该分支初始化正常流率为
\[
\boxed{
f_{p\to q}^0
=
\rho_{p\to q}b_{pY}u_p^0.
}
\tag{31}
\]

固定的产出系数和固定配给比例只缩放物料数量。
当缺口按各自正常流率归一化成“等效正常生产时间”后，
不会改变 \(D(t)\) 的形状；
它们通过正常流率进入库存 coverage
\[
q(0)/f^0.
\]

\section{当前规则的适用边界}

以上闭式与算子形式基于以下 baseline：
\begin{itemize}
    \item \(t<0\) 时网络处于稳态初始化；
    \item production delay 与 transport delay 固定；
    \item 固定生产配方和固定输出配给比例；
    \item 输出边 capacity 不绑定；
    \item 不主动补库、不超额生产、不重新路由；
    \item 正常流恢复到已耗空库存时允许直接穿过，不要求先补回 \(q_m(0)\)；
    \item 没有损耗、报废、保质期约束；
    \item 当前主要考虑完全停产---完全恢复的二元生产强度。
\end{itemize}

若以后允许 partial recovery，
则仍可使用累计缺口函数，但一般有
\[
0\le \dot D(t)\le1,
\]
并且
\[
u_p(t)=u_p^0[1-\dot D_p^{act}(t)]
\]
可以描述部分生产。

若设备故障会直接冻结已经进入流程的 WIP，
则 production delay 不再是 simple pure delay，
需要显式引入 WIP 状态；本文不处理该情形。

\section{STEP1 最终接口}

单窗口下，优先使用直观的对称公式
\[
\boxed{
\begin{aligned}
t_{p\to q}^{\mathrm{stop}}
&=
\sigma_p+\Theta_{p\to q}+C_{p\to q},
\\
t_{p\to q}^{\mathrm{rec}}
&=
\rho_p+\Theta_{p\to q},
\end{aligned}
}
\]
并检查
\[
\rho_p-\sigma_p>C_{p\to q}.
\]

多窗口或一般 DAG 下，使用统一算子
\[
\boxed{
D_{p\to q}(t)
=
\left[
D_p^{act}(t-\Theta_{p\to q})
-C_{p\to q}
\right]_+.
}
\]

因此 STEP1 可以概括为
\[
\boxed{
\underbrace{D_p^{act}}_{\text{上游自身 activity}}
\quad
\xrightarrow{
\substack{
\text{production/transport delay}\\
\text{inventory absorption}
}}
\quad
\underbrace{D_{p\to q}}_{\text{具体下游可见 shortage}}.
}
\]

该 \(D_{p\to q}(t)\) 正是 STEP2 在 required-input merge 时直接接收的输入。

\end{document}
